\section{Vaught's Formula-Algebra Construction of Countable Models}
\label{sec:p03-n-a-003-01}

Vaught's original account begins with Boolean algebras of formulas and
their dual prime ideals. It constructs models that simultaneously omit
specified types and control which formulas their finite tuples satisfy.
The elementary embedding arguments then identify prime and atomic models,
and the saturation construction treats a countable family of relative
interpretations at once \cite[Secs.~1--4, pp.~304--314]{Vaught1961}.
This presentation keeps those constructions and their original
countability conditions explicit.

\subsection{The source's language and formula algebras}
\label{sec:p03-vaught-formula-algebras}

\begin{hbconvention}{def:p03-vaught-standing}{Standing theory and meaning of countability}
Unless a result explicitly changes the setting, \(T\) is a complete
satisfiable theory in an at most countable language \(L\) of finite-arity
relations and individual constants, with logical identity, and has an
infinite model. Completeness then implies that every model is infinite.
Vaught's denumerable models are countably infinite.
The variables are \(v_0,v_1,\ldots\); all substitutions avoid capture.
The constructions use the ambient Choice assumption of
\cref{thm:fol-compactness}.
For theories in these constructions, consistency is satisfiability;
first-order completeness supplies the equivalent syntactic condition.
Countable enumeration here is a set-theoretic enumeration; no algorithm
for deciding consistency or finding isolators is asserted.
The source prints \(\bigvee v_i\) for existential and
\(\bigwedge v_i\) for universal quantification; the displays here
write \(\exists v_i\) and \(\forall v_i\), with the same binding
and truth clauses. These quantifier signs do not denote infinitary
Boolean disjunction or conjunction.
\end{hbconvention}

\begin{hbdefinition}{def:p03-vaught-formula-algebra}{Formulas, validity, atoms and atomistic algebras}
Let \(F_n(T)\) be the formulas with free variables among
\(v_0,\ldots,v_{n-1}\). In this source presentation
\(\vdash_T\varphi\) means validity of the universal closure in every
model of \(T\). Write \(\varphi\approx_T\psi\) when
\(\vdash_T(\varphi\leftrightarrow\psi)\).
The quotient \(\mathbb B_n(T)=F_n(T)/{\approx_T}\), with conjunction,
disjunction and negation, is a Boolean algebra.
We also use the source's \(F_n(T)\) for the corresponding formula
prealgebra, retaining its equivalence relation.

A formula is nonzero, or consistent with \(T\), when
\(T\cup\{\exists v_0\cdots\exists v_{n-1}\varphi\}\) is satisfiable.
An \emph{atom} is a consistent \(\alpha\in F_n(T)\) such that,
for each \(\varphi\in F_n(T)\), consistency of
\(\alpha\land\varphi\) implies \(\vdash_T\alpha\to\varphi\).
An \emph{atomless element} is a consistent \(\beta\) with no atom
\(\alpha\) satisfying \(\vdash_T\alpha\to\beta\).
The algebra is \emph{atomistic} if every consistent formula has an
atom below it. These properties refer to the quotient order; they
do not say that a particular tuple realizes a type.
\end{hbdefinition}

\begin{hbdefinition}{def:p03-vaught-prime-ideals}{The source's dual prime ideals and realized tuples}
A dual prime ideal \(P\) is a nonempty proper subset of \(F_n(T)\),
closed under conjunction and upward under \(T\)-implication, that
contains either \(\varphi\) or \(\neg\varphi\) for every \(\varphi\).
Exactly one of those alternatives belongs to \(P\).
The source calls this a prime ideal; its direction is the one usually
called an ultrafilter. Let \(\mathcal P_n(T)\) be the set of these \(P\).
It is principal if
\[
 P=\{\varphi:\vdash_T\alpha\to\varphi\}
\]
for some consistent \(\alpha\). For \(\mathcal M\models T\), put
\[
 P(\mathcal M)=\{\bar a\in(\dom\mathcal M)^n:
                         \mathcal M\models\varphi(\bar a)
                         \text{ for every }\varphi\in P\}.
\]
Repeated coordinates are allowed.

With distinct new constants \(c_0,\ldots,c_{n-1}\), write
\(\overline\varphi=\varphi(c_0,\ldots,c_{n-1})\) and
\(\overline P=T\cup\{\overline\varphi:\varphi\in P\}\).
The latter is a complete theory in the language with those named
coordinates.
\end{hbdefinition}

\begin{hblemma}{thm:p03-vaught-algebra-types}{proved=annex:p3.vaught.algebra-types}{Atoms, complete types and named expansions}
The dual prime ideals are exactly the consistent complete \(n\)-types
over the empty set, identifying \(\mathcal P_n(T)\) with \(S_n(T)\).
Their source principality agrees with \cref{def:principal-type}.
A principal complete type contains an atom and is isolated by it.
For a \emph{consistent} \(\alpha\), being an atom is equivalent to all
its realizations yielding elementarily equivalent named expansions.
The basic sets \([\varphi]=\{P:\varphi\in P\}\) form the Stone topology;
atomisticity is equivalent to density of isolated types.
\(F_0(T)\) has exactly two equivalence classes and one complete type.
\end{hblemma}

\subsection{The simultaneous Henkin and omission construction}
\label{sec:p03-vaught-simultaneous}

\begin{hbdefinition}{def:p03-vaught-henkin-data}{New constants and witness requirements}
Let \(C=\{c_i:i<\omega\}\) be distinct symbols outside \(L\).
At a finite stage let \(\Gamma_s\) be the conjunction of all accumulated
sentences, consistent with the starting theory.
Enumerate every formula \(\varphi_s(v_0)\) of \(L(C)\) with at most
the displayed free variable. Choose \(c_{\nu_s}\) with least index
absent from \(\varphi_s,\Gamma_s\), and use
\[
 W_s=(\exists v_0\varphi_s)\to\varphi_s(c_{\nu_s}).
\]
Different constants may denote the same element.
An existential block of length zero makes no quantification.
\end{hbdefinition}

\begin{hblemma}{thm:p03-vaught-henkin-naming}{proved=annex:p3.vaught.henkin-naming}{Fresh witnesses, projections and a fully named model}
Adding \(W_s\) preserves consistency at the finite stage.
Replacing distinct new constants by free variables, after freshening
binders, equates consistency of their sentence instance with consistency
of its existential closure.
Any consistent theory in \(L(C)\) containing witnesses for every such
formula has a model whose \(L\)-domain consists of the \(C\)-values.
The named part is elementary in a model of that witness theory.
Thus every consistent countable relational/constant theory has an
at most countable model, and an infinite one when its axioms force
infinitude.
In the fresh-witness construction each \(C\)-value has infinitely
many distinct constant symbols naming it.
\end{hblemma}

\begin{hbtheorem}{thm:p03-vaught-simultaneous-construction}{proved=annex:p3.vaught.simultaneous}{The full original simultaneous construction}
Let \((P_j)_{j\in J}\), \(J\subseteq\omega\), be any specified finite
or countably infinite family of nonprincipal complete types
\(P_j\in\mathcal P_{r_j}(T)\), with finite arities.
The source's \(p_j+1\) is this total arity \(r_j\).
There is a countably infinite \(\mathcal M\models T\) such that:
\begin{enumerate}
\item Every finite tuple satisfies either an atom or an atomless formula
      of its corresponding \(F_n(T)\).
\item \(P_j(\mathcal M)=\varnothing\) for every \(j\in J\).
\end{enumerate}
Either condition may be imposed alone. Arity zero needs no omission
requirement, since its sole complete type is principal.
This is the simultaneous content of source Theorem~2.1, not merely
the existence of a model omitting one type.
\end{hbtheorem}

\begin{hbmethod}{def:p03-vaught-two-branch-construction}{The source's two branches at one stage}
Enumerate every requirement \(\Pi_s=(P,d_0,\ldots,d_{r-1})\), where the
\(d_i\) are distinct \(C\)-symbols and either \(P=\varnothing\) or
\(P\) is a listed type of arity \(r\geq1\).
Pair this enumeration with the witness enumeration.
Let \(\gamma_s=\Gamma_s\land W_s\).
Write \(\bar d\) for the task tuple and \(\bar e\) for the other
occurring \(C\)-symbols, disjoint from \(\bar d\).
Replace these two tuples by free slots \(\bar v,\bar w\), after
freshening binders, obtaining \(\theta_s(\bar v,\bar w)\) in \(L\).
Put \(\beta_s(\bar v)=\exists\bar w\,\theta_s(\bar v,\bar w)\).
It is consistent.

If an atom \(\alpha\) lies below \(\beta_s\), add
\(\alpha(\bar d)\land\gamma_s\).
If \(\beta_s\) is atomless and \(P=\varnothing\), add \(\gamma_s\).
If \(\beta_s\) is atomless and \(P\ne\varnothing\), choose
\(\delta\notin P\) with \(\beta_s\land\delta\) consistent and add
\(\delta(\bar d)\land W_s\), retaining all previous sentences.
The annex proves that this choice always exists and that each branch
preserves consistency. The infinitely many aliases, rather than a
restriction to tuples of distinct elements, allow these syntactically
distinct constant requirements to cover every tuple.
\end{hbmethod}

\begin{hbcorollary}{thm:p03-vaught-automatic-omission}{proved=annex:p3.vaught.automatic-omission}{Further omitted types in the closure of isolated points}
A model with the first property of the simultaneous theorem omits
every nonprincipal complete type whose every neighborhood contains
an isolated point. In this countable formula language these are exactly
the nonprincipal points which are limits of a sequence of distinct
isolated types. This automatically omitted collection may be uncountable;
the separately specified omission list remains countable.
\end{hbcorollary}

\subsection{Elementary maps, prime models and atomic models}
\label{sec:p03-vaught-prime-atomic}

\begin{hbdefinition}{def:p03-vaught-model-notions}{Atomicity, primeness, universality and homogeneity}
An \(\mathcal A\models T\) is atomic when each finite tuple satisfies
an atom of its formula algebra. It is prime when it elementarily embeds
in every model of \(T\). These agree with
\cref{def:p03-atomic-prime} by the algebra/type identification.
The elementary-embedding condition is Vaught's stated prime notion;
his introductory footnote~3 distinguishes earlier ordinary-embedding
variants.
It is \(\aleph_0\)-universal, in the source's convention, when it is
countably infinite and every countably infinite \(T\)-model
elementarily embeds in it.

It is homogeneous when any two finite tuples with the same complete
type are carried one to the other by an automorphism.
Here the data are full first-order types, not just finite relational
isomorphism patterns. An elementary tuple map preserves the truth of
every formula, including identity, with those tuple entries as parameters.
Repeated entries therefore must correspond to repeated entries.
\end{hbdefinition}

\begin{hblemma}{thm:p03-vaught-prefix-embedding}{proved=annex:p3.vaught.prefix-embedding}{The source's complete-prefix embedding criterion}
Suppose \(\dom\mathcal A=\{a_i:i<\omega\}\), allowing repetitions,
and \(\bar b=(b_i)_{i<\omega}\) is in \(\mathcal B\models T\).
If, for every \(n\), \((a_0,\ldots,a_{n-1})\) and
\((b_0,\ldots,b_{n-1})\) have the same complete type, then
\(a_i\mapsto b_i\) is well defined and injective, respects every
constant and relation, and is an elementary embedding.
\end{hblemma}

\begin{hblemma}{thm:p03-vaught-atomic-parameters}{proved=annex:p3.vaught.atomic-parameters}{Naming a finite tuple preserves atomicity}
If \(\mathcal A\) is atomic and \(\bar a\) is a finite tuple from it,
the expansion \((\mathcal A,\bar a)\), considered as a model of its
complete named theory, is atomic. Isolation in that expansion is
over the specified finite parameters; it is not isolation over
an arbitrary infinite parameter set.
\end{hblemma}

\begin{hblemma}{thm:p03-vaught-atomic-embedding}{proved=annex:p3.vaught.atomic-embedding}{Original Lemma 3.1 with its prescribed tuple}
Let \(\mathcal A\) be countably infinite and atomic, and let
\(\mathcal B\) be any \(T\)-model, of any cardinality.
There is an elementary embedding \(\mathcal A\to\mathcal B\).
If finite tuples \(\bar a,\bar b\) have the same complete type,
the embedding can be chosen to send \(\bar a\) to \(\bar b\).
The construction uses atoms of successive source tuples; it imposes
no atomicity or saturation hypothesis on the target.
\end{hblemma}

\begin{hbtheorem}{thm:p03-vaught-atomic-uniqueness}{proved=annex:p3.vaught.atomic-uniqueness}{Original Theorem 3.2: uniqueness by full back-and-forth}
Two countably infinite atomic models of \(T\) are isomorphic.
An elementary correspondence between prescribed finite tuples
can be extended to such an isomorphism.
\end{hbtheorem}

\begin{hbtheorem}{thm:p03-vaught-atomic-homogeneity}{proved=annex:p3.vaught.atomic-homogeneity}{Original Theorem 3.3: elementary homogeneity}
Every countably infinite atomic model of \(T\) is homogeneous
in the stated complete-type sense.
\end{hbtheorem}

\begin{hbtheorem}{thm:p03-vaught-prime-atomic}{proved=annex:p3.vaught.prime-atomic}{Original Theorem 3.4: prime exactly when countable and atomic}
A model of the standing \(T\) is prime if and only if it is
countably infinite and atomic.
In particular, atomicity of an uncountable model does not imply primeness.
Any two prime models are isomorphic.
\end{hbtheorem}

\begin{hbtheorem}{thm:p03-vaught-isolation-density}{proved=annex:p3.vaught.isolation-density}{Original Theorem 3.5: the full existence equivalence}
The following are equivalent:
\begin{enumerate}
\item \(T\) has a prime model.
\item \(T\) has an atomic model, without a cardinal restriction here.
\item Every \(F_n(T)\) is atomistic.
\item For every \(n<\omega\), isolated types are dense in \(S_n(T)\).
\end{enumerate}
The atomic model in clause 2 can have any size; the existence of a
countably infinite atomic model then follows from the simultaneous
construction. This existence equivalence does not assert that the
original arbitrary-sized atomic model itself is prime.
\end{hbtheorem}

\begin{hblemma}{thm:p03-vaught-nonprincipal-algebra}{proved=annex:p3.vaught.nonprincipal-algebra}{Original fact 3.6 with its complete algebra proof}
The countable Boolean algebra \(\mathbb B_n(T)\) has infinitely many
elements if and only if it has a nonprincipal dual prime ideal.
Equivalently, infinitely many inequivalent \(n\)-variable formulas
force a nonisolated complete \(n\)-type.
\end{hblemma}

\begin{hbtheorem}{thm:p03-vaught-finite-formulas-categoricity}{proved=annex:p3.vaught.finite-formulas-categoricity}{The source section 3 consequence of the construction}
For the standing complete countable-language \(T\), the following
are equivalent:
\begin{enumerate}
\item All countably infinite models of \(T\) are isomorphic.
\item Each \(F_n(T)\) has only finitely many inequivalent formulas.
\item Each \(\mathcal P_n(T)\) is finite.
\item Every complete finite-tuple type is principal.
\end{enumerate}
This is the finite-formula characterization discussed after source
Theorem~3.5; its proof uses the supplied atomic and omission constructions.
The original Ryll-Nardzewski, Engeler and Svenonius presentations
remain distinct from this derivation in Vaught's account.
\end{hbtheorem}

\subsection{The exact saturation and parameter scopes}
\label{sec:p03-vaught-saturation}

\begin{hbdefinition}{def:p03-vaught-saturation}{Weak saturation and the source's one-coordinate extension}
A model is weakly saturated when \(P(\mathcal M)\ne\varnothing\)
for every \(P\in\mathcal P_n(T)\) and every finite \(n\).
It is saturated in Vaught's sense when it is weakly saturated and,
whenever \(P\in\mathcal P_n(T)\), \(Q\in\mathcal P_{n+1}(T)\),
\(P\subseteq Q\) under the inclusion of the first \(n\) variable slots,
and \(\bar a\in P(\mathcal M)\), some \(b\) satisfies
\((\bar a,b)\in Q(\mathcal M)\).
Since \(P\) is complete, \(P\subseteq Q\) means
\(Q\cap F_n(T)=P\).

In modern parameter notation, \(\kappa\)-saturation means realization
of every complete finite-tuple type over each parameter set \(A\)
of cardinality less than the specified infinite \(\kappa\).
For \(\kappa=\aleph_0\), this concerns finite parameter sets.
Homogeneity remains the automorphism property of finite elementary
tuple maps defined above; it is not introduced as a synonym for saturation.
\end{hbdefinition}

\begin{hblemma}{thm:p03-vaught-saturation-parameters}{proved=annex:p3.vaught.saturation-parameters}{The finite-parameter meaning of the source definition}
Vaught saturation is equivalent to realization of every complete
finite-tuple type over every finite parameter set, relative to the
complete theory with those parameters named.
After naming an \(m\)-tuple, its complete \(r\)-types inject into
\(\mathcal P_{m+r}(T)\), extending that tuple's fixed type.
Thus smallness of the empty-parameter type spaces also bounds
the spaces over each finite parameter tuple.
For an infinite countable model this does not imply realization
over arbitrary countable parameter sets: a consistent type asking
for a new element different from every existing element is omitted.
\end{hblemma}

\subsection{Relative interpretations and the full original construction}
\label{sec:p03-vaught-relative-saturation}

\begin{hbdefinition}{def:p03-vaught-relative-interpretation}{The source's definable-domain interpretations}
Here \(T_1\) denotes a countable relational/constant target theory.
The interpretation formulas have no external parameters; fixed
parameters must first be named in that target language.
A relative interpretation \(I\) of \(T\) in \(T_1\) specifies a
formula \(D(x)\) for a nonempty domain, a formula \(\rho_R(\bar x)\)
for each relation of \(L\), and a unique distinguished element
in that domain for each individual constant.
Require \(T_1\models\exists x\,D(x)\).
The distinguished element may be given by a target constant or
by a formula \(K_d(x)\) with
\[
                 T_1\models\exists!x\,(D(x)\land K_d(x)).
\]
Equality is actual identity on the domain, not a quotient relation.
Translate quantifiers by
\[
 (\exists x\varphi)^I=\exists x(D(x)\land\varphi^I),\qquad
 (\forall x\varphi)^I=\forall x(D(x)\to\varphi^I),
\]
relations by their specified formulas, and Boolean operations
compositionally. For a defined distinguished element, use a fresh
existential variable satisfying its unique defining formula in the
atomic translation. This also treats repeated constant occurrences
coherently. The interpretation requirement is that
\(T_1\models\sigma^I\) for every axiom \(\sigma\) of \(T\).

In \(\mathcal B\models T_1\), denote the resulting \(T\)-structure
on \(D(\mathcal B)\) by \(I(\mathcal B)\).
These are the definable-domain relative interpretations of source
section 1, with distinguished-element details written explicitly.
\end{hbdefinition}

\begin{hblemma}{thm:p03-vaught-interpretation-preservation}{proved=annex:p3.vaught.interpretation-preservation}{Truth, elementary maps and unions under interpretation}
For assignments in the interpreted domain,
\[
 I(\mathcal B)\models\varphi(\bar a)
                 \quad\Longleftrightarrow\quad
 \mathcal B\models\varphi^I(\bar a).
\]
If \(\mathcal B\preccurlyeq\mathcal B'\), then
\(I(\mathcal B)\preccurlyeq I(\mathcal B')\).
For an elementary increasing chain, interpretation commutes
with its union, including the defined distinguished elements.
\end{hblemma}

\begin{hblemma}{thm:p03-vaught-interpreted-extension}{proved=annex:p3.vaught.interpreted-extension}{Original Lemma 4.1: one extension with its elementary diagram}
Let \(\mathcal B\models T_1\) be countably infinite,
\(\mathcal A=I(\mathcal B)\), and
\(\bar a\in P(\mathcal A)\) with
\(P\in\mathcal P_n(T)\), \(P\subseteq Q\in\mathcal P_{n+1}(T)\).
There is a countably infinite elementary extension
\(\mathcal B^*\succeq\mathcal B\) such that, in
\(\mathcal A^*=I(\mathcal B^*)\), some \(b\) has
\((\bar a,b)\in Q(\mathcal A^*)\).
The empty-prefix case \(n=0\) is included.
\end{hblemma}

\begin{hbtheorem}{thm:p03-vaught-saturated-existence}{proved=annex:p3.vaught.saturated-existence}{Full original Theorem 4.2 for a family of interpretations}
Suppose every \(\mathcal P_n(T)\) is countable.
Let \(T_1\) be a consistent countable theory and let
\((I_i)_{i\in J}\), for a nonempty finite or countable \(J\), be
relative interpretations of \(T\) in \(T_1\).
There is a countably infinite \(\mathcal B\models T_1\) for which
every \(I_i(\mathcal B)\) is a saturated \(T\)-model.
The construction first fills every extension requirement of one
countable stage, and then repeats this operation over an outer
elementary chain to cover new tuples.
The target theory need not be complete.
With the identity interpretation this supplies a countably infinite
saturated model of \(T\).
\end{hbtheorem}

\subsection{Universality, uniqueness and the smallness criterion}
\label{sec:p03-vaught-universality}

\begin{hblemma}{thm:p03-vaught-saturated-embedding}{proved=annex:p3.vaught.saturated-embedding}{Original Lemma 4.3 with a prescribed finite map}
Let \(\mathcal B\) be countably infinite and saturated, and let
\(\mathcal A\models T\) be any countably infinite model.
There is an elementary embedding \(\mathcal A\to\mathcal B\);
it can extend any prescribed correspondence of finite tuples
with the same complete type.
\end{hblemma}

\begin{hbtheorem}{thm:p03-vaught-saturated-uniqueness}{proved=annex:p3.vaught.saturated-uniqueness}{Original Theorem 4.4: uniqueness of the countable saturated model}
Any two countably infinite saturated models of \(T\) are isomorphic.
An elementary correspondence of finite tuples extends to an isomorphism.
\end{hbtheorem}

\begin{hblemma}{thm:p03-vaught-saturated-homogeneity}{proved=annex:p3.vaught.saturated-homogeneity}{Original Lemma 4.5: homogeneous countable saturated models}
Every countably infinite saturated model is homogeneous
in the stated finite complete-type sense.
\end{hblemma}

\begin{hbtheorem}{thm:p03-vaught-saturation-characterization}{proved=annex:p3.vaught.saturation-characterization}{All directions of original Theorem 4.6}
For a countably infinite \(\mathcal A\models T\), the following are
equivalent:
\begin{enumerate}
\item \(\mathcal A\) is saturated.
\item It is \(\aleph_0\)-universal and homogeneous.
\item It is weakly saturated and homogeneous.
\end{enumerate}
The third clause's homogeneity condition is part of the equivalence;
weak realization over the empty set has not been substituted for
the specified finite-parameter extension property.
\end{hbtheorem}

\begin{hbdefinition}{def:p03-vaught-smallness}{The type-counting condition and its stronger neighbor}
The standing theory is \emph{small} when every \(\mathcal P_n(T)\),
\(n<\omega\), is countable. In contrast, \(\omega\)-stability asks
for only countably many complete one-types over every countable
parameter set in a model. The parameter sets and the conditions
are different; smallness does not assert \(\omega\)-stability.
\end{hbdefinition}

\begin{hblemma}{thm:p03-vaught-stability-boundary}{proved=annex:p3.vaught.stability-boundary}{The stronger parameter condition really implies smallness}
For a complete theory in a countable language, \(\omega\)-stability
implies smallness. The converse fails in DLO, as proved below.
This is the parameter-counting comparison for the modern terminology,
not an added hypothesis on original Theorem~4.7.
\end{hblemma}

\begin{hblemma}{thm:p03-vaught-countable-types-atomistic}{proved=annex:p3.vaught.countable-types-atomistic}{The full Boolean-algebra step used in original Corollary 4.8}
If \(\mathcal P_n(T)\) is countable, every nonzero element of
\(\mathbb B_n(T)\) contains an atom.
Indeed an atomless nonzero element would support a binary splitting
tree and produce \(2^{\aleph_0}\) different complete types.
\end{hblemma}

\begin{hbtheorem}{thm:p03-vaught-smallness-equivalences}{proved=annex:p3.vaught.smallness-equivalences}{All four conditions of original Theorem 4.7}
The following are equivalent:
\begin{enumerate}
\item \(T\) is small.
\item It has a countably infinite saturated model.
\item It has an \(\aleph_0\)-universal model.
\item It has a countably infinite weakly saturated model.
\end{enumerate}
The saturated model, when it exists, is unique up to isomorphism
and homogeneous. Saturation retains the finite-parameter meaning
proved above.
\end{hbtheorem}

\begin{hbcorollary}{thm:p03-vaught-universal-prime}{proved=annex:p3.vaught.universal-prime}{Original Corollary 4.8}
If \(T\) has an \(\aleph_0\)-universal model, then it has a prime model.
Equivalently, smallness implies the prime-existence condition here.
The universal model and the prime model are not required to coincide.
\end{hbcorollary}

\subsection{Worked models and the boundaries of the hypotheses}
\label{sec:p03-vaught-examples}

\begin{hbtheorem}{thm:p03-vaught-dlo-example}{proved=annex:p3.vaught.dlo-example}{DLO, finite patterns, named cuts and the stability distinction}
Let \(T_{\rm dlo}\) be dense linear order without endpoints in the
language consisting of \(<\) and identity.
Each complete \(n\)-type over the empty set is isolated by its tuple's
equality and strict-order pattern. Their number is
\[
              \sum_{k=0}^n k!\,\left\{\begin{matrix}n\\k\end{matrix}\right\},
\]
where the braces count partitions of an \(n\)-element set into
\(k\) nonempty blocks; at \(n=0\) the number is one.
The ordered rational model is prime, atomic, saturated and homogeneous.
It is the unique countably infinite model.

Nevertheless there are \(2^{\aleph_0}\) complete one-types over the
countable rational parameter set: distinct real cuts give distinct types.
Thus \(T_{\rm dlo}\) is small but not \(\omega\)-stable.
In the countable language naming every rational, its complete expanded
theory \(T_{\mathbb Q}\) still has the prime atomic model \(\mathbb Q\),
but has \(2^{\aleph_0}\) complete one-types over the empty set of
\emph{that expanded language}. Each irrational cut type is nonprincipal
and omitted by the rational model. Therefore prime existence does not
imply smallness, and the finite-parameter smallness bridge cannot be
extended to arbitrary countable naming.
\end{hbtheorem}

\begin{hbtheorem}{thm:p03-vaught-discrete-example}{proved=annex:p3.vaught.discrete-example}{A source-native prime model and a nonprincipal extension obstruction}
Let \(T_{\rm disc}=\operatorname{Th}(\omega,<)\).
For each \(k<\omega\), the formula saying that \(x\) has exactly
\(k\) predecessors isolates the type of \(k\).
The countably infinite ordered model \((\omega,<)\) is atomic and prime.
The partial type saying that \(x\) has at least \(k\) predecessors
for every \(k\) is consistent and nonprincipal.
Each complete extension of it is nonprincipal and omitted by
that prime model. Thus a prime model need not be weakly saturated.
This works directly with the completed discrete-order example;
its unique finite-position formulas supply the prime-model calculation
without a decision procedure for the full classification of types.
\end{hbtheorem}

\begin{hbtheorem}{thm:p03-vaught-unary-obstruction}{proved=annex:p3.vaught.unary-obstruction}{A countable language with atomless one-variable formulas}
Let \(T_{\rm bits}\) say that every finite Boolean pattern of the
unary predicates \(U_0,U_1,\ldots\) has infinitely many realizations.
Its language has only these predicates and logical identity.
It is a complete satisfiable theory.
Every consistent one-variable formula is atomless.
There are \(2^{\aleph_0}\) complete one-types, one for each
infinite predicate bit sequence, and no prime or atomic model.
The simultaneous construction's atomless branch is indispensable
in this example.
\end{hbtheorem}

\begin{hbtheorem}{thm:p03-vaught-finite-boundary}{proved=annex:p3.vaught.finite-boundary}{Finite complete theories and the cardinal restriction on atomicity}
If a complete theory \(S\) in a countable relational/constant language
has a finite nonempty model \(\mathcal A\), every \(S\)-model is
isomorphic to \(\mathcal A\). Its model is prime, atomic,
homogeneous and saturated over every parameter set within it.
This is a separate finite case; Vaught's denumerable equivalence
was stated with infinite models.

For the complete theory of infinite pure equality, every model
of every infinite cardinality is atomic, whereas its prime models
are precisely the countably infinite ones.
Completeness also matters: in the pure equality language, for the
incomplete theory saying only
that the universe has one or two elements, the one-element model
is atomic in the corresponding formula-algebra sense but cannot
elementarily embed in a two-element model.
\end{hbtheorem}

The source also illustrates its criteria with discrete orders and
real closed fields \cite[p.~315]{Vaught1961}.
For discrete orders it invokes an outside decision procedure for
the full type count and reports the universal homogeneous order
\(\omega+(\omega^*+\omega)\cdot\eta\), where \(\eta\) is the
countable dense order type without endpoints.
For real closed fields it cites the real algebraic field as prime
while the one-type space has continuum many members.
Those additional application claims use the corresponding
order-decision and real-closed-field results. The direct examples
above distinguish prime existence, isolation and smallness.

\subsection{Functional signatures and the broader omission interface}
\label{sec:p03-vaught-language-bridge}

\begin{hblemma}{thm:p03-vaught-functional-bridge}{proved=annex:p3.vaught.functional-bridge}{A complete bridge to countable languages with total functions}
Replace each positive-arity function \(f\) by a relation \(G_f\)
for its graph, with the total and single-valued graph axioms;
zero-arity functions remain individual constants.
There is a cardinality-preserving bijection between the function-language
models of a theory and their graph-language presentations.
Formula translation preserves truth, consistency and free-variable arity,
and has an inverse modulo the graph axioms.
It identifies complete and partial types, principality and isolation,
elementary maps, primeness, atomicity, saturation and homogeneity.
The graph language is countable when the original one is.
Consequently the preceding results transport to countable
function languages with these proved identifications.
The original source presentation remains relational/constant.
\end{hblemma}

\begin{hbtheorem}{thm:p03-vaught-partial-omission}{proved=annex:p3.vaught.partial-omission}{The general countable partial-type omission construction}
Let \(S\) be any satisfiable theory in a countable first-order language,
not necessarily complete, and let \((p_j)_{j\in J}\) be a finite or
countable family of nonprincipal partial types, consistent with \(S\),
in finite tuples of variables.
Then \(S\) has an at most countable model omitting all of them.
It is infinite when \(S\) forces infinitude.
Here principality has the exact consistent-formula meaning of
\cref{def:principal-type}, not merely nonisolation of one selected completion.
Zero-arity types are included.
This generalization uses the witness and finite-refinement construction;
it is distinct from the source's complete-theory simultaneous atom/atomless
theorem and supplies the broader interface of \cref{thm:omitting-types}.
\end{hbtheorem}
